\chapter*{Preface}
\addcontentsline{toc}{chapter}{Preface}

\emph{A preface...}

\begin{flushright}
{\makeatletter\itshape
    \@author \\
    Capital, \monthname{} \the\year{}
\makeatother}
\end{flushright}

My initial reason for choosing engineering was my belief that it was a field built around scientific research, practical experimentation, and critical thinking. However, my experience at university was somewhat shocking. In many cases, engineering education felt like an extension of high school: memorize the equations, memorize the professor’s approach, and reproduce it in the exam. As a result, summer training represents a welcome opportunity for me to learn, conduct independent research, and genuinely understand the subjects I study without constantly worrying about whether I am following an academic approach that I find ineffective.

I am not particularly interested in cloud computing itself, but I believe that learning a new field is valuable, especially when it introduces me to technologies and concepts outside my usual interests. Since this training covers Linux, Python, virtualization, Docker, Kubernetes, networking, and cloud computing, I decided to take part in it and broaden my technical background.

My main interests, however, lie elsewhere. I am particularly interested in PCB design, embedded systems, simulation, robotics, and systems planning. I also enjoy designing simple mechanical components and enclosures for my PCBs using FreeCAD. Outside engineering, I enjoy Linux ricing and often spend considerable time customizing and experimenting with my system. In my free time, I also enjoy playing Minecraft.

There is another subject that concerns me deeply: the widespread availability of artificial intelligence. Although I rely heavily on AI in my daily life, I am uncomfortable with the idea of making it universally available without meaningful restrictions. I believe that excessive dependence on AI can make people less capable of thinking independently, less willing to struggle through difficult problems, and potentially less creative. More importantly, I worry about what happens when we begin delegating not only our repetitive tasks, but also our thinking, creativity, and decision-making to machines.

Sometimes I wonder what purpose remains for a person when an artificial entity can surpass them in the very things they value most about themselves. I expect that, if current trends continue, AI may eventually become responsible for an increasing portion of human intellectual work. In such a future, there is a risk that many people will simply become spectators of achievements they no longer meaningfully participate in creating.

At times, I see myself as a trivial person fighting windmills alone. At other times, I believe that I am not. Perhaps neither judgment is meaningful at my age. I still have a great deal to learn, and I have not yet reached a position from which I can meaningfully influence the world around me. What matters to me is reaching that stage first. I would like, eventually, to contribute to the discussion surrounding how powerful technologies such as AI should be developed, distributed, and used. In particular, I believe that access to such technologies should involve responsibility, experience, and a genuine reason for using them rather than unrestricted availability for everyone.

Finally, I would like to thank my family and friends for their support throughout this difficult period. Their support has made this journey considerably easier, and I am genuinely grateful for it.

